\documentclass{uflamon}          % classe base para a monografia

%==============================================================================
% Utilizacao de pacotes
\usepackage[T1]{fontenc}         % usa fontes postscript com acentos
\usepackage[brazil]{babel}       % hifenização e títulos em português do Brasil
\usepackage[utf8]{inputenc}     % permite edição direta com acentos
\usepackage{amsmath}             % pacote da AMS para Matemática Avançada
\usepackage{amssymb}             % símbolos extras da AMS
\usepackage{latexsym}            % símbolos extras do LaTeX
\usepackage{graphicx}            % para inserção de gráficos
\usepackage{listings}            % para inserção de código
\usepackage{fancyvrb}            % para inserção de saídas de comandos
%\usepackage{enumerate}           % para personalizar lista enumeradas 
											%(incluso na classe)
\usepackage{longtable}           % para tambelas muito grandes NOVO!!!!

\usepackage{colortbl} % cores em tabelas
\newcolumntype{Z}{|>{\columncolor[gray]{0.9}}l|} %cor cinza em células
%\usepackage{array} % já incluso na classe
\newcolumntype{L}[1]{>{\raggedright\let\newline\\\arraybackslash\hspace{0pt}}m{#1}}
\newcolumntype{C}[1]{>{\centering\let\newline\\\arraybackslash\hspace{0pt}}m{#1}}
\newcolumntype{R}[1]{>{\raggedleft\let\newline\\\arraybackslash\hspace{0pt}}m{#1}}
\usepackage{multirow} % para juntar duas linhas em uma só
\usepackage{multicol} % para uso de várias colunas
% \usepackage{caption}

% cores para os links cruzados
\usepackage{color}
\definecolor{rltred}{rgb}{0.2,0,0}
\definecolor{rltgreen}{rgb}{0,0.2,0}
\definecolor{rltblue}{rgb}{0,0,0.2}

\usepackage[colorlinks=true,
            urlcolor=rltblue,       % \href{...}{...} external (URL)
            filecolor=rltgreen,     % \href{...} local file
            linkcolor=black,       % \ref{...} and \pageref{...}
            citecolor=black,
            pdftitle={Filtro Prensa},
          pdfauthor={Marcus Bruno Fernandes Silva},
          pdfsubject={Planta Piloto de Filtro Prensa},
          pdfkeywords={filtro prensa}%
]{hyperref} % para referência cruzadas
%\usepackage{hyperref}            % para referência cruzadas
% \usepackage{subfigure}           % figuras dentro de figuras
\usepackage{caption}            % remodelando o formato dos títulos de 
\usepackage{subcaption}            % entra em conflito com subfigure

% configuração padrão do listings   
\lstset{%
   language=Java,
   extendedchars=true,
   tabsize=3,
   basicstyle=\footnotesize\ttfamily,
   stringstyle=\em,
   showstringspaces=false 
}

%abnt-emphasize=bf coloca o título das bibliografias em negrito
%abnt-thesis-year=both
\usepackage[alf,num,abnt-etal-cite=3,abnt-etal-list=0,abnt-url-package=url,abnt-emphasize=bf]{abntex2cite}

% redefinindo formatação de títulos de tabelas e figuras

%==============================================================================
\usepackage{microtype} 			% para melhorias de justificação
\usepackage{float}
\usepackage{titling}
\usepackage[version=4]{mhchem}
\usepackage{array}
\usepackage{units}
\usepackage{siunitx}
\usepackage{enumitem}
\usepackage{booktabs}
% \usepackage{nameref} % Referencia capitulos por nome
%\usepackage[authoryear]{natbib}

% Especificando hifenizações que por ventura LaTeX não saiba fazer
% Por padrão 99,9% dos termos em português devem ser hifenizados corretamente.
\hyphenation{hardware software Li-nux am-bien-te diag-nos-ti-car coor-de-na-ção 
  FAE-PE Recovery TelEduc Williams UFLA Sérgio Renata Reynolds disciplina Iara
Hernandez Rodrigues}

%==============================================================================
% Dados da monografia, capa: autor, titulo, banca, etc... - SUBSTITUA DE ACORDO
%==============================================================================
\author{Gabriela Nunes \\
        Gustavo Henrique de Souza Paiva \\
        Karen Leticia Sanchez Costa \\
        Larissa Souza Castro \\
        Marcus Bruno Fernandes Silva \\
        Marina Leticia Alves Ferreira}
\subtitle{Exemplo para os Usuários}
%\author{Joao}
\title{Planta Piloto de Filtro Prensa}
\engtitle{Use of Uflamon Class}
%\engsubtitle{Sample for Users}
%\edicao{3$^a$ edição revista, atualizada e ampliada}
\date{2018}
\newcommand{\orient}{Professora Zuy Maria Magriotis}
\newcommand{\disciplina}{GNE347}
\tipo{Pŕe-projeto entregue e apresentado à \orient{} como requisito da disciplina \disciplina{}.}
\orientador{\orient}
\newcommand{\otipo}{Pŕe-projeto entregue e apresentado à \orient{} como requisito da disciplina \disciplina{}.}
\newcommand{\olocal}{Lavras -- MG}
\local{Lavras -- MG}

%##################################################

\newcommand \fazercapa{
\begin{titlepage}
    \centering

    \includegraphics[scale=0.45]{logoufla.jpg}%

    \vspace{1cm}

    \Large\theauthor
    
    \vspace{\stretch{1}}

    \bfseries\LARGE\MakeUppercase\thetitle

    \vspace{\stretch{1}}

    \Large\MakeUppercase\olocal\\
    \MakeUppercase\thedate
    \null
    \cleardoublepage
\end{titlepage}
}

\newcommand \fazerfolhaderosto{
    \begin{titlepage}
        \newpage

        \begin{center}
            \MakeUppercase\theauthor

            \vspace{\stretch{1}}

            \textbf{\MakeUppercase\thetitle}

            \vspace{\stretch{1}}
            
            \hfill\begin{minipage}{8cm}{\otipo}\end{minipage} % tipo da monografia

            \vspace{\stretch{.58}}

            \ifthenelse{\equal{\@orientador}{\@empty}} % % se orientador nao existir,
            {}%                                            % nao produza nada
            {\orient\\
             Orientadora
            }

            \vspace{\stretch{.86}} 

            \textbf{\MakeUppercase\olocal}\\ 
            \textbf{\MakeUppercase\thedate}
            \end{center}
    \end{titlepage}
}

\newcolumntype{P}[1]{>{\centering\arraybackslash}p{#1}}
\graphicspath{{./figuras/}}

\newcommand{\tabitem}{~~\llap{\textbullet}~~}
\newcommand{\R}{$^{\tiny{\textregistered}}$~}


%==============================================================================

% Meus comandos

\newcommand{\citeay}[1]{\citeauthoronline{#1} (\citeyear{#1})}
\newcommand{\citeayp}[1]{(\citeauthoronline{#1}, \citeyear{#1})}
\newcommand{\citean}[1]{\citeauthoronline{#1} \cite{#1}}
\newcommand{\pderr}[2]{\frac{\partial}{\partial #2}(#1)}
\newcommand{\pa}[1]{\left(#1 \right)}

%==============================================================================
% Começar o documento

%\newsubfloat{figure}

%\includeonly{capitulos/materiais_metodos,capitulos/resultados}
\begin{document}

\fazercapa

\fazerfolhaderosto

\tableofcontents

\clearpage

\pagestyle{ufla}

%==============================================================================
% incluindo os capitulos


\chapter{Introdução} \label{sec:intro}

\section{Referencial Teórico}\label{sec:refteo}


%%% Local Variables:
%%% mode: latex
%%% TeX-master: "../main_archive"
%%% End:

\chapter{Motivação}\label{motivacao}


%%% Local Variables:
%%% mode: latex
%%% TeX-master: "../main_archive"
%%% End:

\chapter{Descrição do Projeto}\label{descricao}

\section{Materiais utilizados}

% Please add the following required packages to your document preamble:
% \usepackage{booktabs}
\begin{table}[H]
\centering
\caption{Materiais requeridos para o projeto}
\label{tabela:custo}
\begin{tabular}{@{}lr@{}}
\toprule
\multicolumn{1}{c}{\textbf{Materiais}}  & \multicolumn{1}{c}{\textbf{Custo (R\$)}} \\ \midrule
Tubos PVC ($\SI{50}{\mm}$)              & 7,00/m                                   \\
Joelho $\ang{90}$ ($\SI{50}{\mm}$)      & 3,00/unidade                             \\
Tubos de ensaio                         & 5,00/unidade                             \\
Garras                                  & 25,00/unidade                            \\
Suportes universais                     & 50,00/unidade                            \\
Papel milimetrado                       & 0,50/folha                               \\
Tela de tecido                          & 4,00/m                                   \\
Termômetro                              & 130,00/unidade                           \\
Relógio                                 & -                                        \\
Sílica gel                              & 29,00/kg                                 \\
Água                                    & -                                        \\
Acetona                                 & 40,00/L                                  \\
Etanol                                  & 22,00/L                                  \\ \bottomrule
\end{tabular}
\end{table}


\section{Procedimento experimental}

Inicialmente o tubo de PVC deverá ser seccionado, por uma furadeira, no formato
circular dos tubos de ensaio que serão usados para armazenar as soluções
analisadas.

O tubo então deverá ser fixado em hastes metálicas, e nas entradas de ar do tubo
deverão ser instalados telas contendo sílica, a fim de evitar a passagem de
umidade externa pelo tubo.

Pretende-se realizar testes com pelo menos três tipos de substâncias dispostas
simultaneamente em tubos de ensaios com escala milimetrada. Serão testadas a
influência do tipo de convecção natural e forçada (com o auxílio de um secador
de cabelos), e a influência da temperatura (os tubos de ensaio deverão estar
submersos em banho maria a temperatura constante).


%%% Local Variables:
%%% mode: latex
%%% TeX-master: "../main_archive"
%%% End:

\chapter{Procedimento Experimental}\label{procedimento}



%%% Local Variables:
%%% mode: latex
%%% TeX-master: "../main_archive"
%%% End:

\chapter{Resultados e Discussões}
\label{chap:resultados}


%%% Local Variables:
%%% mode: latex
%%% TeX-master: "../main_archive"
%%% End:

\chapter{Conclusão}
\label{chap:conclusao}


%%% Local Variables:
%%% mode: latex
%%% TeX-master: "../main_archive"
%%% End:


%==============================================================================

% Incluindo bibliografia
%\bibliographystyle{plain}             % estilo para labels em numeros
%\bibliographystyle{alpha}             % estilo para labels em iniciais
\bibliographystyle{abntex2-num}         %% estilo para referências usando ABNT, 
 %para usar o abntex2-alf precisa arrumar o nome dos autores na bibliografia
%\bibliographystyle{abntex2-alf}           % estilo para referências usando ABNT, 

%inclui Referências Bibliográficas
\referencias
\bibliography{bibliografia/b}

%==============================================================================
% Fim do texto
\end{document}

%%% Local Variables:
%%% mode: latex
%%% TeX-master: t
%%% End:
